\begin{table*}[t]
\centering
\tablefont
\caption{\ct{Baseline Patient, Clinical, Procedural, and Outcome Characteristics Across the National STS ACSD v2.81 Cohorts (1,091 U.S.\ Surgical Centers, 2014--2017).}}
\label{tab:cohort_baseline}
\setlength{\tabcolsep}{5.0pt}
\renewcommand{\arraystretch}{1.04}
\begin{tabularx}{\textwidth}{@{}Yrrrr@{}}
\toprule
\thd{Characteristic} & \thd{Full Development} & \thd{Nested 50k Benchmark} & \thd{Unseen External Test} & \thd{SMD$^{\dagger}$} \\
 & \thd{($N=346{,}454$; 1,017 ctr)} & \thd{($N=50{,}000$; 1,017 ctr)} & \thd{($N=21{,}241$; 74 ctr)} & \thd{(Dev vs.\ Test)} \\
\midrule
\multicolumn{5}{@{}l}{\textit{Demographics and Anthropometry}} \\
\quad Age, mean $\pm$ SD, yr & 65.1 $\pm$ 10.1 & 66.4 $\pm$ 10.0 & 65.2 $\pm$ 10.1 & 0.009 \\
\quad Female sex & 87,423 (25.2\%) & 12,093 (24.2\%) & 5,172 (24.4\%) & 0.020 \\
\quad White race & 285,838 (82.5\%) & 41,977 (84.0\%) & 17,206 (81.0\%) & 0.039 \\
\quad Black race & 26,750 (7.7\%) & 3,413 (6.8\%) & 1,822 (8.6\%) & 0.031 \\
\quad Asian race & 11,889 (3.4\%) & 1,648 (3.3\%) & 842 (4.0\%) & 0.028 \\
\quad Hispanic or Latino ethnicity & 26,921 (7.8\%) & 3,583 (7.2\%) & 1,433 (6.8\%) & 0.039 \\
\quad Body-mass index, mean $\pm$ SD, kg/m$^2$ & 30.1 $\pm$ 6.0 & 30.1 $\pm$ 6.0 & 30.0 $\pm$ 6.0 & 0.014 \\
\quad Body surface area, mean $\pm$ SD, m$^2$ & 2.01 $\pm$ 0.24 & 2.02 $\pm$ 0.24 & 2.01 $\pm$ 0.24 & 0.008 \\
\addlinespace[3pt]
\multicolumn{5}{@{}l}{\textit{Preoperative Clinical Comorbidities and Cardiac Substrate}} \\
\quad Hypertension & 306,810 (88.6\%) & 44,569 (89.1\%) & 19,032 (89.6\%) & 0.033 \\
\quad Diabetes mellitus & 168,522 (48.6\%) & 24,025 (48.0\%) & 10,456 (49.2\%) & 0.012 \\
\quad Heart failure (current or prior) & 75,070 (21.7\%) & 11,333 (22.7\%) & 5,018 (23.6\%) & 0.047 \\
\quad NYHA class III or IV & 35,307 (10.2\%) & 5,449 (10.9\%) & 2,511 (11.8\%) & 0.052 \\
\quad Prior myocardial infarction & 180,339 (52.0\%) & 26,134 (52.3\%) & 10,960 (51.6\%) & 0.009 \\
\quad Chronic lung disease & 82,357 (23.8\%) & 12,200 (24.4\%) & 4,643 (21.9\%) & 0.045 \\
\quad Peripheral vascular disease & 46,703 (13.5\%) & 7,052 (14.1\%) & 2,688 (12.7\%) & 0.025 \\
\quad Cerebrovascular disease & 68,057 (19.6\%) & 10,323 (20.6\%) & 3,959 (18.6\%) & 0.025 \\
\quad Chronic kidney disease (eGFR $<60$) & 72,267 (20.9\%) & 11,068 (22.1\%) & 4,561 (21.5\%) & 0.015 \\
\quad Preoperative dialysis & 9,939 (2.9\%) & 1,435 (2.9\%) & 612 (2.9\%) & 0.001 \\
\quad Recorded preoperative AF history & 1,762 (0.5\%) & 282 (0.6\%) & 138 (0.7\%) & 0.018 \\
\quad Left ventricular EF, mean $\pm$ SD, \% & 52.4 $\pm$ 12.2 & 52.3 $\pm$ 12.3 & 52.4 $\pm$ 12.1 & 0.001 \\
\quad Serum creatinine, mean $\pm$ SD, mg/dL & 1.18 $\pm$ 1.06 & 1.19 $\pm$ 1.06 & 1.20 $\pm$ 1.10 & 0.018 \\
\quad CKD-EPI 2021 eGFR, mL/min/1.73\,m$^2$ & 77.7 $\pm$ 23.5 & 76.5 $\pm$ 23.4 & 77.1 $\pm$ 23.6 & 0.025 \\
\addlinespace[3pt]
\multicolumn{5}{@{}l}{\textit{Surgical Acuity, Anatomy, and Preoperative Medications}} \\
\quad Three-vessel coronary disease & 264,116 (76.2\%) & 38,363 (76.7\%) & 16,465 (77.5\%) & 0.031 \\
\quad Elective operative status & 129,234 (37.3\%) & 18,690 (37.4\%) & 7,887 (37.1\%) & 0.004 \\
\quad Urgent operative status & 202,366 (58.4\%) & 29,130 (58.3\%) & 12,439 (58.6\%) & 0.003 \\
\quad Emergent or salvage status & 14,810 (4.3\%) & 2,173 (4.3\%) & 910 (4.3\%) & 0.001 \\
\quad Preoperative $\beta$-blocker therapy & 312,486 (90.2\%) & 45,027 (90.0\%) & 18,507 (87.1\%) & 0.097 \\
\quad Preoperative amiodarone therapy & 21,722 (6.3\%) & 2,935 (5.9\%) & 1,120 (5.3\%) & 0.043 \\
\addlinespace[3pt]
\multicolumn{5}{@{}l}{\textit{Primary Outcome (Treatment-Requiring New-Onset POAF)}} \\
\quad \textbf{POAF events (STS \texttt{COtAFib})} & \textbf{84,177 (24.30\%)} & \textbf{25,000 (50.0\%)} & \textbf{5,160 (24.29\%)} & \textbf{0.0002} \\
\bottomrule
\end{tabularx}
\par\smallskip\raggedright\fontsize{7.5}{9.4}\selectfont\color{sub}
$^{\dagger}$Standardized mean difference (SMD) comparing the natural-prior full development cohort (1,017 hospitals) with the strictly hospital-disjoint unseen external validation cohort (74 hospitals; 0.0\% hospital overlap). Every clinical characteristic exhibits $|\mathrm{SMD}| < 0.10$, and all nine U.S.\ Census divisions match within 0.18 percentage points. The nested 50,000-patient benchmark is class-balanced (25,000 POAF / 25,000 controls) strictly within the same 1,017 development hospitals and recalibrated back to natural clinical prevalence using the development-only Bayes offset ($-1.1365$).
\end{table*}
